\subsection{Software architectuur}
Om de onderhoudbaarheid te verbeteren van het huidige CMS is er ook nagedacht over het opzetten van de code.
De principe houd het volgende in: 
\begin{enumerate}
    \item \textbf{Single-responsibility principle}
    Dit betekent dat een class, maar 1 verantwoordelijk mag hebben en bij definitie ook maar 1 taak.
    Als een class te veel verantwoordelijkheden heeft dan wordt het moeilijker om de code te begrijpen en aan te passen.

    \item \textbf{Open-closed principle}
    Dit houdt in als class of functie of een andere software uitgebreid moet worden dat het wordt gedaan door middel van een extentie(open) in plaats van modificatie(closed).
    Hierdoor hou je oude code intact en heb je geen risico dat je bestaande code stuk gaat vanwege de nieuwe functionaliteit die geschreven is. 
        
    \item \textbf{Liskov substitution principle} 
    Een class die afgeleid is van een base class, zou moeten vervangen kunnen worden voor  een andere instantie van die base class.
    Dit moet gedaan kunnen worden zonder de validiteit (corectness) van het programma te beïnvloeden.
    Door het gebruik van LSP verhoog je de consistentiteit en de verwachte uitkomst van je progamma.

    \item \textbf{Interface segregation} 
    Een interface moet alleen de methods geven die nodig is voor de client. 
    Geen client moet geforceerd zijn om methodes te implementeren waar die geen gebruik van kan maken.
    Dit uit zich vaak in plaats van grote interface gebruik meerdere kleinere interfaces.
    Deze interfaces zijn verantwoordelijk voor meer specifike usecases inplaats van generaliseerde usecases

    \item \textbf{Dependency inversion} 
    Een class of module zou niet moeten afhangen van implementaties maar van abstracties.
    Hier door verminder je de koppeling van de modules/classes, en verhoog je de code onderhoudbaarheid.
    Om aan dit principe te voldoen wordt er vaak gebruik gemaakt van dependency injection.

\end{enumerate}

\whitespace[2]
Deze principes hebben geleid tot de keuze van een handler software architectuur.
De handler structuur heeft elementen van het Chain of responsibility design pattern.
Hierbij heeft elke handler 1 functie waar die aan voldoet.
Hierdoor blijft de code makkelijk te begrijpen en beter onderhoudbaar.

\begin{graphic}
    \captionsetup{type=figure}
    \caption{Sequence diagram Handler structuur}
    \includegraphics[scale=0.2]{Placeholder.jpg}
    \label{fig:SequenceDiagramHandlerStructure}
\end{graphic}
\whitespace
- misschien software entiteit gebruiken plaats van class
- solid
- handlers
- units of work aan kaarten
plaatje als voor beeld

